\documentclass[11pt]{article}

\usepackage[utf8]{inputenc}
\usepackage[T1]{fontenc}
\usepackage[spanish]{babel}
\usepackage{amsmath,amssymb,mathtools,bm,graphicx,xcolor}
\usepackage{float}
\usepackage{svg}
\svgsetup{inkscapelatex=false}
\usepackage[a4paper,margin=2.5cm]{geometry}
\usepackage[
    colorlinks=true,
    linkcolor=red,
    citecolor=green,
    urlcolor=red
]{hyperref}

\spanishdecimal{.}

\setlength{\parindent}{0pt}
\setlength{\parskip}{0.45em}
\setlength{\abovedisplayskip}{6pt}
\setlength{\belowdisplayskip}{6pt}
\setlength{\abovedisplayshortskip}{4pt}
\setlength{\belowdisplayshortskip}{4pt}

% Formato de encabezado y problemas
\newcommand{\encabezado}[3]{%
{\Large\bfseries #1\par}
\vspace{2pt}
{\normalsize #2\hfill #3\par}
\vspace{6pt}\hrule\vspace{8pt}}

\newcommand{\problema}[2]{%
\vspace{0.55em}
\noindent\textbf{#1.}\enspace{\small #2}\par
\vspace{0.3em}}

\newcommand{\inciso}[1]{%
\vspace{0.35em}
\noindent\textbf{(#1)}\enspace}

% Comandos auxiliares
\newcommand{\dd}{\,\mathrm{d}}
\newcommand{\ii}{\mathrm{i}}
\newcommand{\e}{\mathrm{e}}
\newcommand{\sen}{\operatorname{sen}}
\newcommand{\grad}{\nabla}
\newcommand{\laplaciano}{\nabla^2}
\newcommand{\R}{\mathbb{R}}
\newcommand{\N}{\mathbb{N}}
\newcommand{\Nzero}{\mathbb{N}_0}
\newcommand{\JR}{J.R.}
\newcommand{\GR}{G.R.}

\begin{document}

\encabezado{Difusión del calor 2D en un rectángulo}{Física Matemática II}{J.R., G.R.}

\problema{1}{Planteamiento del problema}

Consideremos la ecuación de difusión del calor bidimensional en el dominio rectangular
\begin{equation*}
\Omega=\{(x,y)\in\R^2:0<x<a,\;0<y<b\},
    \qquad a>0,\quad b>0,
\end{equation*}
para tiempos $t\geq 0$. La función desconocida será
\begin{equation*}
\psi=\psi(x,y,t),
\end{equation*}
la cual representa la temperatura del sistema en el punto $(x,y)$ y en el instante $t$.

Siguiendo la notación usual del apunte del curso, escribimos la ecuación del calor en la forma
\begin{equation}
\frac{\partial^2\psi}{\partial x^2}
    +
    \frac{\partial^2\psi}{\partial y^2}
    -
    \frac{1}{\alpha}\frac{\partial\psi}{\partial t}=0,
    \qquad \alpha>0,
    \label{eq:calor_2d_rubilar}
\end{equation}
o, equivalentemente,
\begin{equation*}
\frac{\partial\psi}{\partial t}
    =
    \alpha\left(
    \frac{\partial^2\psi}{\partial x^2}
    +
    \frac{\partial^2\psi}{\partial y^2}
    \right).
\end{equation*}
Aquí $\alpha$ corresponde a la difusividad térmica del material.

Impondremos condiciones de borde homogéneas tipo Dirichlet en todo el contorno de $\Omega$:
\begin{align}
\psi(0,y,t)&=0,
    &\psi(a,y,t)&=0,
    &&0\leq y\leq b,
    \quad t\geq0,
    \label{eq:cb_x_calor}\\
    \psi(x,0,t)&=0,
    &\psi(x,b,t)&=0,
    &&0\leq x\leq a,
    \quad t\geq0.
    \label{eq:cb_y_calor}
\end{align}
La condición inicial se escribe como
\begin{equation*}
\psi(\vec{x},0)=\varphi(\vec{x}),
    \qquad \vec{x}=(x,y),
\end{equation*}
o, de forma equivalente,
\begin{equation}
\psi(x,y,0)=\varphi(x,y).
    \label{eq:ci_calor_xy}
\end{equation}

\problema{2}{Aplicación del método de separación de variables}

Buscamos una solución separable de la forma
\begin{equation}
\psi_{\mathrm{sep}}(x,y,t)=X(x)Y(y)T(t),
    \label{eq:ansatz_calor}
\end{equation}
donde $X$, $Y$ y $T$ dependen sólo de una variable. Sustituimos \eqref{eq:ansatz_calor} en la ecuación \eqref{eq:calor_2d_rubilar}. Primero calculamos
\begin{equation*}
\frac{\partial^2\psi_{\mathrm{sep}}}{\partial x^2}=X''(x)Y(y)T(t),
    \qquad
    \frac{\partial^2\psi_{\mathrm{sep}}}{\partial y^2}=X(x)Y''(y)T(t),
\end{equation*}
y
\begin{equation*}
\frac{\partial\psi_{\mathrm{sep}}}{\partial t}=X(x)Y(y)T'(t).
\end{equation*}
Por lo tanto, al reemplazar en \eqref{eq:calor_2d_rubilar}, obtenemos
\begin{equation*}
X''YT+XY''T-\frac{1}{\alpha}XYT'=0.
\end{equation*}
Dividiendo por $\psi_{\mathrm{sep}}=XYT$, suponiendo que trabajamos donde esta cantidad no se anula idénticamente, resulta
\begin{equation*}
\frac{X''}{X}+\frac{Y''}{Y}-\frac{1}{\alpha}\frac{T'}{T}=0.
\end{equation*}
Luego,
\begin{equation*}
\frac{1}{\alpha}\frac{T'}{T}
    =
    \frac{X''}{X}+\frac{Y''}{Y}.
\end{equation*}
El miembro izquierdo depende sólo de $t$, mientras que el miembro derecho depende sólo de las variables espaciales $x$ e $y$. Como la igualdad debe valer para todo $(x,y,t)$ del dominio, ambos miembros deben ser iguales a una constante de separación. Para que la parte temporal sea decreciente, escribimos
\begin{equation*}
\frac{1}{\alpha}\frac{T'}{T}=-\lambda,
    \qquad
    \frac{X''}{X}+\frac{Y''}{Y}=-\lambda.
\end{equation*}
De la primera ecuación se obtiene inmediatamente
\begin{equation*}
T'+\alpha\lambda T=0.
\end{equation*}

La parte espacial todavía contiene dos variables. Para separarla, escribimos
\begin{equation*}
\frac{X''}{X}+\lambda=-\frac{Y''}{Y}.
\end{equation*}
El lado izquierdo depende sólo de $x$ y el derecho sólo de $y$. Por lo tanto, ambos deben ser iguales a una nueva constante de separación. Es conveniente escribir
\begin{equation*}
\frac{X''}{X}=-\lambda_x,
    \qquad
    \frac{Y''}{Y}=-\lambda_y,
    \qquad
    \lambda=\lambda_x+\lambda_y.
\end{equation*}
Así se obtiene el sistema de EDOs
\begin{align}
X''+\lambda_x X&=0,
    \notag\\
    Y''+\lambda_y Y&=0,
    \notag\\
    T'+\alpha(\lambda_x+\lambda_y)T&=0.\label{eq:edo_t_calor}
\end{align}

\problema{3}{Condiciones de borde homogéneas y autovalores}

Las condiciones de borde homogéneas pueden imponerse directamente sobre las soluciones separables. En efecto, usando \eqref{eq:ansatz_calor}, las condiciones \eqref{eq:cb_x_calor} implican
\begin{equation*}
X(0)Y(y)T(t)=0,
    \qquad
    X(a)Y(y)T(t)=0.
\end{equation*}
Para soluciones no triviales, $Y\not\equiv0$ y $T\not\equiv0$, de modo que
\begin{equation*}
X(0)=0,
    \qquad
    X(a)=0.
\end{equation*}
Análogamente, de \eqref{eq:cb_y_calor} se obtiene
\begin{equation*}
Y(0)=0,
    \qquad
    Y(b)=0.
\end{equation*}

Por lo tanto, debemos resolver los problemas de autovalores
\begin{equation*}
X''+\lambda_xX=0,
    \qquad X(0)=0,
    \qquad X(a)=0,
\end{equation*}
y
\begin{equation*}
Y''+\lambda_yY=0,
    \qquad Y(0)=0,
    \qquad Y(b)=0.
\end{equation*}

\inciso{a}\textbf{Dirección $x$}\par

Analizamos el signo de $\lambda_x$, pues de ello depende la forma explícita de la solución.

Si $\lambda_x=0$, entonces
\begin{equation*}
X''=0
    \quad\Longrightarrow\quad
    X(x)=A+Bx.
\end{equation*}
Las condiciones $X(0)=0$ y $X(a)=0$ implican $A=0$ y $B=0$, por lo que sólo se obtiene la solución trivial.

Si $\lambda_x<0$, escribimos $\lambda_x=-p^2$, con $p>0$. Entonces
\begin{equation*}
X''-p^2X=0
    \quad\Longrightarrow\quad
    X(x)=A\e^{px}+B\e^{-px}.
\end{equation*}
La condición $X(0)=0$ entrega $A+B=0$. Luego,
\begin{equation*}
X(a)=A(\e^{pa}-\e^{-pa})=0.
\end{equation*}
Como $\e^{pa}-\e^{-pa}\neq0$, se concluye que $A=0$ y $B=0$. Nuevamente sólo aparece la solución trivial.

Por lo tanto, el caso no trivial requiere $\lambda_x>0$. En tal caso,
\begin{equation*}
X(x)=A\cos(\sqrt{\lambda_x}\,x)+B\sin(\sqrt{\lambda_x}\,x).
\end{equation*}
De $X(0)=0$ se obtiene $A=0$. Luego, la condición $X(a)=0$ exige
\begin{equation*}
B\sin(\sqrt{\lambda_x}\,a)=0.
\end{equation*}
Para no obtener $B=0$, debe cumplirse
\begin{equation*}
\sqrt{\lambda_x}\,a=n\pi,
    \qquad n=1,2,3,\ldots.
\end{equation*}
Así, los autovalores y autofunciones en la dirección $x$ son
\begin{equation*}
\lambda_{x,n}=\left(\frac{n\pi}{a}\right)^2,
    \qquad
    X_n(x)=\sin\left(\frac{n\pi x}{a}\right),
    \qquad n=1,2,3,\ldots.
\end{equation*}

\inciso{b}\textbf{Dirección $y$}\par

El cálculo para $Y(y)$ es completamente análogo. Las condiciones $Y(0)=0$ y $Y(b)=0$ seleccionan
\begin{equation*}
\lambda_{y,m}=\left(\frac{m\pi}{b}\right)^2,
    \qquad
    Y_m(y)=\sin\left(\frac{m\pi y}{b}\right),
    \qquad m=1,2,3,\ldots.
\end{equation*}
Por lo tanto, cada modo espacial queda etiquetado por un par $(n,m)\in\N^2$ y posee autovalor total
\begin{equation}
\lambda_{nm}=\lambda_{x,n}+\lambda_{y,m}
    =
    \pi^2\left(\frac{n^2}{a^2}+\frac{m^2}{b^2}\right).
    \label{eq:lambda_nm_calor}
\end{equation}

\problema{4}{Parte temporal}

Sustituyendo \eqref{eq:lambda_nm_calor} en \eqref{eq:edo_t_calor}, se obtiene
\begin{equation*}
\frac{\dd T_{nm}}{\dd t}
    +
    \alpha\pi^2\left(\frac{n^2}{a^2}+\frac{m^2}{b^2}\right)T_{nm}=0.
\end{equation*}
Esta ecuación diferencial ordinaria lineal de primer orden se resuelve directamente:
\begin{equation*}
T_{nm}(t)=C_{nm}\exp\left[-\alpha\pi^2
    \left(\frac{n^2}{a^2}+\frac{m^2}{b^2}\right)t\right].
\end{equation*}

Con esto, una solución separable asociada al par $(n,m)$ es
\begin{equation*}
\psi^{\mathrm{sep}}_{nm}(x,y,t)
    =
    C_{nm}
    \sin\left(\frac{n\pi x}{a}\right)
    \sin\left(\frac{m\pi y}{b}\right)
    \exp\left[-\alpha\pi^2
    \left(\frac{n^2}{a^2}+\frac{m^2}{b^2}\right)t\right].
\end{equation*}

\problema{5}{Superposición general y condición inicial}

Como la ecuación \eqref{eq:calor_2d_rubilar} es lineal y homogénea, una combinación lineal de soluciones separables también es solución. Entonces escribimos
\begin{equation}
\psi(x,y,t)
    =
    \sum_{n=1}^{\infty}\sum_{m=1}^{\infty}
    A_{nm}
    \sin\left(\frac{n\pi x}{a}\right)
    \sin\left(\frac{m\pi y}{b}\right)
    \exp\left[-\alpha\pi^2
    \left(\frac{n^2}{a^2}+\frac{m^2}{b^2}\right)t\right].
    \label{eq:solucion_general_calor}
\end{equation}
Esta expresión ya satisface la EDP y todas las condiciones de borde homogéneas. Falta imponer la condición inicial \eqref{eq:ci_calor_xy}. Evaluando en $t=0$,
\begin{equation}
\varphi(x,y)
    =
    \sum_{n=1}^{\infty}\sum_{m=1}^{\infty}
    A_{nm}
    \sin\left(\frac{n\pi x}{a}\right)
    \sin\left(\frac{m\pi y}{b}\right).
    \label{eq:serie_seno_doble}
\end{equation}
Por lo tanto, los coeficientes $A_{nm}$ son los coeficientes de Fourier seno dobles de la condición inicial $\varphi$.

Para calcularlos, multiplicamos \eqref{eq:serie_seno_doble} por
\begin{equation*}
\sin\left(\frac{n'\pi x}{a}\right)
    \sin\left(\frac{m'\pi y}{b}\right)
\end{equation*}
e integramos sobre el rectángulo $[0,a]\times[0,b]$:
\begin{align*}
&\int_0^a\int_0^b
    \varphi(x,y)
    \sin\left(\frac{n'\pi x}{a}\right)
    \sin\left(\frac{m'\pi y}{b}\right)
    \dd y\dd x
    \notag\\
    &=
    \sum_{n=1}^{\infty}\sum_{m=1}^{\infty}A_{nm}
    \left[\int_0^a
    \sin\left(\frac{n\pi x}{a}\right)
    \sin\left(\frac{n'\pi x}{a}\right)\dd x\right]
    \notag\\
    &\hspace{8em}\times
    \left[\int_0^b
    \sin\left(\frac{m\pi y}{b}\right)
    \sin\left(\frac{m'\pi y}{b}\right)\dd y\right].
\end{align*}
Usando las relaciones de ortogonalidad
\begin{equation*}
\int_0^a
    \sin\left(\frac{n\pi x}{a}\right)
    \sin\left(\frac{n'\pi x}{a}\right)\dd x
    =\frac{a}{2}\delta_{nn'},
\end{equation*}
\begin{equation*}
\int_0^b
    \sin\left(\frac{m\pi y}{b}\right)
    \sin\left(\frac{m'\pi y}{b}\right)\dd y
    =\frac{b}{2}\delta_{mm'},
\end{equation*}
obtenemos finalmente
\begin{equation}
A_{nm}
    =
    \frac{4}{ab}
    \int_0^a\int_0^b
    \varphi(x,y)
    \sin\left(\frac{n\pi x}{a}\right)
    \sin\left(\frac{m\pi y}{b}\right)
    \dd y\dd x.
    \label{eq:A_nm_calor}
\end{equation}

Con \eqref{eq:solucion_general_calor} y \eqref{eq:A_nm_calor} queda determinada la solución del problema.


\problema{6}{Gráficos}

\begin{figure}[H]
    \centering
    \includesvg[width=0.62\linewidth]{figuras/calor_2d_modo_11}
    \caption{Primer modo espacial permitido por las condiciones homogéneas de Dirichlet. Se grafica el producto $X_1(x)Y_1(y)=\sen(\pi x/a)\sen(\pi y/b)$ para $a=1$ y $b=2$. La figura permite verificar visualmente que el modo se anula en los cuatro lados del rectángulo: $x=0$, $x=a$, $y=0$ e $y=b$. Como corresponde al modo $(n,m)=(1,1)$, no aparecen nodos interiores; por eso representa la variación espacial más suave compatible con el borde fijo a temperatura nula.}
    \label{fig:calor-modo-11}
\end{figure}

\begin{figure}[H]
    \centering
    \includesvg[width=0.98\linewidth]{figuras/calor_2d_evolucion_seno_gaussiana}
    \caption{Evolución temporal de una condición inicial suave y localizada, evaluada mediante la solución en serie doble. Los parámetros son $a=1$, $b=2$, $\alpha=1$, $T_0=1$, $x_0=0.6a$, $y_0=0.7b$, $\sigma_x=0.12a$ y $\sigma_y=0.13b$. El punto negro marca el centro inicial de la perturbación térmica. La escala de color se mantiene fija en todos los paneles; por tanto, la disminución de intensidad no es un artefacto de la escala, sino el decaimiento real de la temperatura. La difusión ensancha el perfil y reduce su amplitud, mientras las condiciones de borde mantienen el valor nulo en el contorno.}
    \label{fig:calor-evolucion}
\end{figure}

\begin{figure}[H]
    \centering
    \includesvg[width=0.62\linewidth]{figuras/calor_2d_tasas_decaimiento}
    \caption{Tasas de decaimiento $\gamma_{nm}=\alpha\pi^2(n^2/a^2+m^2/b^2)$ para los primeros modos de la serie, usando $a=1$, $b=2$ y $\alpha=1$. Cada casilla muestra el valor numérico de $\gamma_{nm}$. Los modos con índices mayores decaen más rápido porque tienen mayor curvatura espacial. Además, como $b>a$, una variación modal en la dirección $y$ contribuye menos a la tasa que una variación equivalente en la dirección $x$.}
    \label{fig:calor-tasas}
\end{figure}

\end{document}
